\documentclass{ximera}
\title{The Chain Rule and the Price of a Stock}

\newcommand{\pskip}{\vskip 0.1 in}

\begin{document}
\begin{abstract}
A stock problem with the chain rule.
\end{abstract}
\maketitle


\begin{question} \label{QLfel33mz}
The function
\[
     P = f(t) = 64 - 20\cos\left(  \frac{1}{4} t \right) \, , \, 8\leq t \leq 16,
\]
expresses the price of a stock (in dollars/share) in terms of the number of hours past 9am.

\begin{enumerate}

\item What are the units of the constant $64$ in the function above? The constant $20$? The constant $1/4$? Explain your reasoning.

\item Sketch by hand a reasonably accurate graph of the function. Label the axes with the appropriate variable names and units. Include scales on both axes. Label the coordinates of any turning points.

\item Find an expression for the derivative $dP/dt$. Do this by making an explicit $u$-substitution and showing each step of the chain rule. But do \emph{not} use $u$. Use a more meaningful variable instead.

\item What are the units of the derivative $dP/dt$? Interpret its meaning.

\item Find the stock price when it is decreasing at the rate of ($\$3$/share)/hr.

\item Find a function 
\[
   P = g(r)
\]
that expresses the stock price (in dollars/share) in terms of the rate of change in the price (measured in (dollars/share)/hr). Simplify the expression for your function.

\item Sketch by hand a graph of the function $P=g(r)$. Label the axes with the appropriate variable names and units. 

\item Evaluate the derivative 
\[
  \frac{dP}{dr}\Big|_{r=-3}.
\] 
Include units and interpret its meaning as a scaling factor. Include an example of how the scaling factor acts on a \emph{specific} small change.

%\item Is the stock price increasing or decreasing the second time the stock sells for $\$45$/share? At what rate?

\item Is the number of shares you can buy with $\$1000$ increasing or decreasing when the price is decreasing at the rate of ($\$4$/share)/hour? 

\begin{enumerate}

\item Use numerical methods, no calculus, to approximate the rate of change (with respect to time) in the number of shares you can buy with $\$1,000$ when the price is decreasing at the rate of ($\$4$/share)/hour. No calculus.

\item Use calculus to compute the rate in part (i). Quotient rule not allowed.

\end{enumerate}
\end{enumerate}

\begin{explanation}
\begin{enumerate}
\item Because $P$ has units dollars/share, so does the constant $64$. And because the outputs to the cosine function are dimensionless, the coefficient $20$ has these same units.

Because the inputs to the cosine function are measured in radians (ie. are dimensionless) and $t$ is measured in hours, the units 
of the coefficient $1/4$ are rad/hour or $\text{hr}^{-1}$.

\item The function describes an oscillation about the mean price of $64$ dollars/share, with a maximum oscillation about the mean of $20$ dollars/share and a period of 
\[
  \frac{2\pi \text{rad}}{\frac{1}{4}\text{ rad/hr}} = 8\pi \text{ hrs}.
\]

\item The derivative (omitting some of the required steps) is
\begin{align*}
 \frac{dP}{dt} &= \frac{d}{dt} \left(  64 - 20\cos\left(  \frac{1}{4} t \right)\right) \\
                    &= 20\sin \left(  \frac{1}{4} t \right) \frac{d}{dt}\left( \frac{1}{4}t  \right) \\
                    &= 5 \sin \left(  \frac{1}{4} t \right) .
\end{align*}

\item The units of $dP/dt$ are (dollars/share)/hour. The derivative measures the rate at which the stock price is changing with respect to time. Multiplying the derivative by a small time interval (in hours) gives a good approximation to the change in price (in dollars/sh) over that time interval.

\item Suppose the stock price is decreasing at the rate of 3 (dollars/share)/hr at time $t=t_0$ hours past 9am.
Then
\begin{align*}
  -3  &= \frac{dP}{dt}\Big|_{t=t_0} \\
       &=  5 \sin \left(  \frac{1}{4} t \right)\Big|_{t=t_0} \\
                                         &= 5 \sin \left(  \frac{1}{4} t_0 \right) .
\end{align*}
So
\[
           \sin \left(  \frac{1}{4} t_0 \right) = -\frac{3}{5}
\]
and
\begin{align*}
  \cos \left(  \frac{1}{4} t_0 \right) &= \pm \sqrt{1 - \sin^2 \left(  \frac{1}{4} t_0 \right)} \\
                                                  &= \pm \frac{4}{5} .
\end{align*}

To choose the correct sign, we need to remember the domain of the function $P=f(t)$. Since 
\[
    8 \leq t_0 \leq 16 ,
\]
\[
  2 \leq \frac{t_0}{4} \leq 4 .
\]
So
\[
     \frac{\pi}{2} < \frac{t_0}{4} \leq  \frac{3\pi}{2}
\]
and therefore
\[
  \cos \left(  \frac{1}{4} t_0 \right) < 0.
\]
So
\[
    \cos \left(  \frac{1}{4} t_0 \right) = - \frac{4}{5}.
\] 
Finally, the stock price at this time is
\begin{align*}
    f(t_0) &= 64 - 20\cos\left(  \frac{1}{4} t_0 \right) \\
             &= 64 + 16 \\
             &= 80.
\end{align*}
So when the stock price is decreasing at the rate of $3$ (dollars/share)/hour, the stock is selling at a price of $80$ dollars/share.

\item The solution to this part is much the same as the solution to part (f). We just replace $-3/5$ (dollars/share)/hr with $r$ (dollars/share)/hr).

Omitting the details , we get
\[
      \cos (t_0) = -\sqrt{1-\frac{r^2}{25}} = -\frac{1}{5}\sqrt{25 - r^2},
\]
where $t_0$ is the time (measured in hours past noon) when the price is increasing at the rate of $r$ (dollars/share)/hr. 

And the stock price at this time is
\[
    P = f(t_0) = 64 + 4\sqrt{25-r^2}.
\]
and our function is
\[
   P = g(r) = 64 + 4\sqrt{25-r^2}.
\]

\item The graph is part of a semi-circle opening downward. I'll leave it to you to find the correct domain of this function.

\item 
\begin{align*}
  \frac{dP}{dr}\Big|_{r=-3} &=  \frac{d}{dr}\left( 64 + 4\sqrt{25-r^2} \right) \Big|_{r=-3} \\
                                        &=- \frac{4r}{\sqrt{25 - r^2}}\Big|_{r=-3} \\
                                        &= 3.
\end{align*} 
The derivative is 
\[
   3 \frac{\text{dollars/share}}{\text{(dollars/share)/hour}} .
\]

{\bf LEAVE THE UNITS LIKE THIS. DO NOT SIMPLIFY THEM.}

So, for example, during the time it takes the rate of change in the stock price to increase from $-3$ (dollars/share)/hr to $-2.90$ (dollars/share)/hr, the price of the stock increases by approximately $0.30$ dollars/share.  


\end{enumerate}
\end{explanation}

\end{question}

\end{document}

